%==============================================================================
%  Short-form unbalanced-pooling appendix.  Included directly from main-sections.tex.
%
%  Pulls auto-generated rank-condition shares (\unbUrbanRate*, \unbResidShare*,
%  \unbShare*) from RP7/scripts/1b_unbalanced_rank_diagnostic.do.
%==============================================================================

\section{Consistency under unbalanced pooling}\label{app:pooling}
\setcounter{assumption}{0}

% Auto-generated by RP7/scripts/1b_unbalanced_rank_diagnostic.do
% Do not edit by hand; rerun the do-file to refresh.
%
% Rank-condition diagnostics for the unbalanced pooling step
% (Proposition~\ref{prop:pooling}, Appendix~\ref{app:pooling}).
% See quality_reports/2026-05-02_rank-condition-diagnostic.md.
\newcommand{\unbCountCHN}{20,532}
\newcommand{\unbBothCHN}{1,884}
\newcommand{\unbShareCHN}{9.2}
\newcommand{\unbAlwaysRuralCHN}{9,658}
\newcommand{\unbAlwaysUrbanCHN}{8,990}
\newcommand{\unbUrbanRateCHN}{48.0}
\newcommand{\unbResidShareCHN}{90.6}
\newcommand{\unbCountIDN}{26,432}
\newcommand{\unbBothIDN}{6,917}
\newcommand{\unbShareIDN}{26.2}
\newcommand{\unbAlwaysRuralIDN}{10,249}
\newcommand{\unbAlwaysUrbanIDN}{9,266}
\newcommand{\unbUrbanRateIDN}{47.9}
\newcommand{\unbResidShareIDN}{90.2}
\newcommand{\unbCountTZA}{3,170}
\newcommand{\unbBothTZA}{182}
\newcommand{\unbShareTZA}{5.7}
\newcommand{\unbAlwaysRuralTZA}{1,968}
\newcommand{\unbAlwaysUrbanTZA}{1,020}
\newcommand{\unbUrbanRateTZA}{35.0}
\newcommand{\unbResidShareTZA}{88.6}


Many individuals in the country panels are observed in strictly fewer than $T_c$ waves.
This appendix formalizes how we incorporate them into the restricted GRC of Section~\ref{subsec:restricted-grc-model}.
The trajectories $\underline d_i \in \mathcal D$ are only defined for individuals who we observe in all $T_c$ waves, so we augment Equation~\eqref{eq:restricted-GRC} with a single dummy for this group:
\begin{equation}
\label{eq:restricted-grc-unbalanced}
\begin{aligned}
y_{it} =\; &
\sum_{\underline d \in \mathcal D \setminus \{d_T\}}
    \mu_{\underline d}\,\mathbbm 1\{\underline d_i=\underline d\}
+ \Delta_{\underline d_0}\,D_{it} \\
& + \sum_{\underline d \in \mathcal D_S \setminus \{\underline d_0\}}
    \phi\,(\mu_{\underline d}-\mu_{\underline d_0})\,D_{it}\,
    \mathbbm 1\{\underline d_i=\underline d\}
+ \kappa_T\,D_{it}\,
    \mathbbm 1\{\underline d_i=d_T\} \\
& + \mu_{\mathrm{unb}}\,U_i
+ \big(\Delta_{\mathrm{unb}}-\Delta_{\underline d_0}\big)\,U_i\,D_{it}
+ x_{it}'\gamma + \varepsilon_{it},
\end{aligned}
\end{equation}
where $U_i = 1$ if $i$ is observed in strictly fewer than $T_c$ waves and $U_i = 0$ otherwise, $\mu_{\mathrm{unb}}$ and $\Delta_{\mathrm{unb}}$ are the average rural consumption and return to urban location for the unbalanced individuals, respectively, and $\kappa_T \equiv \mu_{d_T} + \phi(\mu_{d_T}-\mu_{\underline d_0})$ is the composite always-urban coefficient from Equation~\eqref{eq:restricted-GRC}.
On the balanced subsample, Equation~\eqref{eq:restricted-grc-unbalanced} collapses to the restricted GRC of Section~\ref{subsec:restricted-grc-model}.
In the unbalanced subsample, the trajectory indicators drop out and the equation collapses to $y_{it} = \mu_{\mathrm{unb}} + \Delta_{\mathrm{unb}} D_{it} + x_{it}'\gamma + \varepsilon_{it}$.

The pooled specification relies on two conditions: first, we need urban location to be exogenous given covariates within the unbalanced cell:

\begin{assumption}[Observed-period orthogonality]
\label{ass:orthogonality}
At the true parameter,
$E[\varepsilon_{it}(\theta_0)\mid U_i=1,\,D_{it},\,x_{it}] = 0$.
\end{assumption}

Since the pooled-cell parameters $(\mu_{\mathrm{unb}},\Delta_{\mathrm{unb}})$ absorb group-level differences in mean rural consumption and mean returns to urban between the two strata, the assumption only restricts the residual conditional mean.
The assumption holds if comparative advantage among the unbalanced does not vary systematically with urban location, once we condition on $x_{it}$.
% This is a stronger requirement than the trajectory-cell version of the assumption on the balanced stratum, where $\theta_i$ is absorbed by the cell.

Second, Equation~\eqref{eq:restricted-grc-unbalanced} encodes a common-$\gamma$ restriction: the same coefficient vector $\gamma$ governs the partial effect of $x_{it}$ on the balanced and unbalanced strata.
The restriction lets us estimate $\gamma$ jointly from both subsamples instead of separately on each.
It holds when the partial effect of $x_{it}$ on log consumption does not vary with attrition status---for instance, when individuals who attrit realize the same returns to education as those who stay.
% The condition is testable: replace $x_{it}'\gamma$ with $x_{it}'\gamma + U_i\,x_{it}'\delta$ and test $\delta = 0$.

The argument follows \citet{verdierAverageTreatmentEffects2020a}, who develops GMM estimation of correlated random coefficient models on unbalanced panels in his Appendix~G.
Verdier uses an individual-fixed-effects parameterization; we use a pooled cell so that $\{\Delta_{\underline d}\}_{\underline d \in \mathcal D_S}$ remains a direct estimation output of \eqref{eq:restricted-grc-unbalanced}.
Beyond that change of bookkeeping, the moment system aggregates within person and across individuals as in Verdier.
Assumption~\ref{ass:orthogonality} together with the common-$\gamma$ restriction deliver moment validity on both strata, and standard non-linear GMM arguments \citep{hansen1982gmm,neweyMcFadden1994} deliver the asymptotic distribution.
The trajectory-block moments identify $(\Delta_{\underline d_0}, \phi, \{\mu_{\underline d}\}_{\underline d \in \mathcal D \setminus \{d_T\}}, \kappa_T)$ as in Section~\ref{subsec:restricted-grc-model}; the $(U_i, U_iD_{it})$ moments identify $(\mu_{\mathrm{unb}}, \Delta_{\mathrm{unb}})$ under the rank condition discussed below; the $x_{it}$ moments identify $\gamma$ under the common-$\gamma$ restriction.
Inference follows the main paper's convention; $V$ is the cluster-robust two-step GMM variance reported there.

\begin{proposition}[Pooling balanced and unbalanced individuals]
\label{prop:pooling}
Under (A1)--(A5) of Section~\ref{subsec:empirical-model}, Assumption~\ref{ass:orthogonality}, the common-$\gamma$ restriction encoded in Equation~\eqref{eq:restricted-grc-unbalanced}, the rank condition that $D_{it}$ has positive residual variance among unbalanced observations ($U_i = 1$) after partialling out $x_{it}$, and standard regularity conditions for non-linear GMM, as $n \to \infty$,
\[
\sqrt n\,(\hat\theta - \theta_0) \xrightarrow{d} \mathcal N(0,V),
\]
where $\theta_0 = (\Delta_{\underline d_0},\phi,\{\mu_{\underline d}\}_{\underline d \in \mathcal D \setminus \{d_T\}},\kappa_T,\mu_{\mathrm{unb}},\Delta_{\mathrm{unb}},\gamma)$ and $V$ is the efficient two-step GMM asymptotic variance.
\end{proposition}

Proposition~\ref{prop:pooling} says that the pooled GMM estimator from \eqref{eq:restricted-grc-unbalanced} is consistent and asymptotically normal for every parameter, including the pooled-cell pair $(\mu_{\mathrm{unb}}, \Delta_{\mathrm{unb}})$.
Two consequences for the parameters of interest follow by the delta method.
The switcher returns $\Delta_{\underline d} = \Delta_{\underline d_0} + \phi(\mu_{\underline d}-\mu_{\underline d_0})$ for $\underline d \in \mathcal D_S$ inherit consistency and asymptotic normality directly.
For the always-urban cell, the LCA restriction $\Delta_{d_T} = \Delta_{\underline d_0} + \phi(\mu_{d_T}-\mu_{\underline d_0})$ combined with the definition of $\kappa_T$ inverts to
\[
\mu_{d_T} = \frac{\kappa_T + \phi\,\mu_{\underline d_0}}{1+\phi}, \qquad \Delta_{d_T} = \Delta_{\underline d_0} + \phi(\mu_{d_T}-\mu_{\underline d_0}),
\]
which is smooth in $(\kappa_T, \phi, \mu_{\underline d_0})$ provided $\phi \neq -1$.
Consistency and asymptotic normality of $(\hat\mu_{d_T}, \hat\Delta_{d_T})$ follow by the delta method on the inversion.

Identification of $\Delta_{\mathrm{unb}}$ requires variation in $D_{it}$ among unbalanced person-periods after partialling out $x_{it}$.
If $D_{it}$ is collinear with $(U_i, x_{it})$ on the unbalanced stratum, the $U_iD_{it}$ moment carries no residual signal and $\Delta_{\mathrm{unb}}$ drops out.
We report two diagnostics on this rank condition.

The marginal share of unbalanced person-periods with $D=1$ is \unbUrbanRateCHN\% in China, \unbUrbanRateIDN\% in Indonesia, and \unbUrbanRateTZA\% in Tanzania, well inside $(0,1)$ in every sample.
With $x_{it}$ matched to the main-estimation covariate vector, $D_{it}$ retains \unbResidShareCHN\% of its within-stratum variation in China, \unbResidShareIDN\% in Indonesia, and \unbResidShareTZA\% in Tanzania after partialling out $x_{it}$.

A complementary diagnostic is the share of unbalanced individuals observed with both $D=0$ and $D=1$ across their own waves.
This share is \unbShareCHN\% in China, \unbShareIDN\% in Indonesia, and \unbShareTZA\% in Tanzania.
The smaller within-switch shares for China and Tanzania mean $\Delta_{\mathrm{unb}}$ is identified primarily off between-individual variation in those samples; in Indonesia about a quarter of unbalanced individuals contribute within-individual variation directly.
